\section*{September 21, 2026: Initial assessed values do not remove the pre-closure dip}
\textit{Drafted by Codex; check requested by Jacob.}

Conditioning on each property's 2006 assessed value does not remove the
relative log-price decline before closure. On identical transactions,
allowing initially different values to predict different annual price paths
reduces the pooled dollar estimate from \$27,769 to \$11,301. The log
estimate changes from 2.7\% to 1.3\%. These results leave the differential
recession/recovery concern unresolved; the adjustment does not identify
the cause of the divergence.

The source is Cook County's Assessor--Assessed Values dataset,
\texttt{uzyt-m557}, extracted September 21, 2026. The control is the log of
\texttt{mailed\_tot} for tax year 2006, joined by the unique 14-digit parcel
PIN. Mailed values precede assessment appeals. They remain assessed dollars,
without conversion to market values. A fixed 2006 value precedes every
sale outcome and the recession. The contemporary historical extract can
include subsequent administrative corrections.

The starting sample remains the 11,202 half-mile single-family and townhouse
sales in 2008--2018. A usable initial house value requires positive total and
building assessments and a single-family property class in 2006. This leaves
1,541 of 1,686 treated sales and 8,911 of 9,516 control sales: 10,452 transactions
at 22 treated and 45 control sites. Of the 750 exclusions, 211 lack a historical
record, 43 lack a positive total and 496 have another historical property class,
including 226 land-class and 218 apartment-building cases. No later value is
substituted. Cleaning, prices, treatment definitions and site assignments in
the working pipeline are unchanged.

\begin{center}
\begin{tabular}{lrrr}
 & Dollar OLS & Log OLS & PPML \\
Full sample, existing model & \$34,851 & 3.0\% & 2.1\% \\
Matched sample, existing model & \$27,769 & 2.7\% & 1.1\% \\
Add initial value, fixed coefficient & \$28,737 & 3.4\% & 1.3\% \\
Add initial value, yearly coefficients & \$11,301 & 1.3\% & 0.02\% \\
All premiums vary by year & \$6,372 & $-1.6\%$ & $-1.2\%$ \\
\end{tabular}
\end{center}

Pooled models compare 2014--2018 with 2008--2012, omitting 2013. All four
matched specifications use exactly 9,462 transactions. Every pooled 95\%
interval includes zero. With yearly initial-value coefficients the dollar
interval is $[-\$18{,}700,\$41{,}302]$ and the log interval is
$[-14.2\%,19.6\%]$. A constant initial-value coefficient adjusts levels;
yearly coefficients additionally allow recovery to vary with initial value.
All models retain site and year effects, transaction weights and site clustering.
The last row also varies every existing characteristic coefficient by year.

The 2010 matched log coefficient changes from $-0.237$ to $-0.220$ after
the yearly assessment adjustment: $-21.1\%$ to $-19.7\%$ relative to 2012,
with an adjusted interval of $[-31.5\%,-6.0\%]$. Allowing all premiums to
vary leaves a $-19.6\%$ gap. The joint log pre-test changes from p=0.025 to
p=0.086; this does not establish parallel trends. The full-premium dollar
model rejects its pre-test (p=0.026). Among treated sales the median fixed
2006 assessment is \$26,345 in the 2012 reference year versus \$17,958 in
2010, showing a change in the homes sold. This adjustment nevertheless
leaves the log dip. Assessments do not measure foreclosure exposure or
subsequent neighborhood shocks.

{\small Reproduction: run Make in
\texttt{tasks/audits/home\_initial\_assessment/code} (working tree based on
\texttt{e76f4dd}). Original source responses, metadata and queries are frozen
in an archive; ordinary builds reuse it. The six-page packet separates
coverage and control changes. Six baseline models replicate, independent OLS
reproduces the new pooled dollar/log and annual log coefficients, and PPML
convergence and score checks pass. Standard reports accompany the saved data.}
